\section{Introducción}

\subsection{Propósito}
Este documento define los requerimientos del prototipo de \textbf{Resilis}, un sistema de apoyo a la decisión para las autoridades y organismos que coordinan la respuesta tras un desastre natural. Sirve como acuerdo de alcance para construir, en un tiempo limitado, un prototipo funcional que demuestre la viabilidad técnica de la propuesta presentada al Rally Latinoamericano de Innovación 2026.

Está pensado para ser leído por el autor del desarrollo y por el jurado de la competencia. Por eso es breve y usa un lenguaje sencillo. La estructura sigue de forma simplificada las normas IEEE 830 e ISO/IEC/IEEE 29148.

\subsection{Alcance del producto}
Resilis ayuda a decidir a quién atender primero, dónde concentrar los esfuerzos y qué recursos asignar durante la post-emergencia, incluso cuando la información es incompleta. Su diferencia frente a otros sistemas es que no solo prioriza: también muestra \emph{qué tan confiable es cada dato}, señala \emph{dónde falta información}, permite \emph{verificarla} y \emph{da seguimiento a cada persona que pide ayuda hasta que la recibe}, con un código de entrega de un solo uso.

La solución se apoya en tres funciones centrales:
\begin{enumerate}
  \item \textbf{Registrar las donaciones} y llevar el inventario de lo que hay disponible, con trazabilidad hasta el destino.
  \item \textbf{Priorizar las necesidades} de barrios y de personas, con su nivel de confianza.
  \item \textbf{Optimizar los recursos humanitarios}: calcular cómo repartir lo disponible entre las necesidades para cubrir la mayor necesidad posible, con la aprobación de la entidad.
\end{enumerate}

El prototipo cubre:
\begin{itemize}
  \item el registro de usuarios identificados, con sus reportes de daños y sus solicitudes de ayuda;
  \item el cálculo de un puntaje de prioridad y un nivel de confianza por barrio y por solicitud;
  \item un mapa de barrios con la prioridad, la confianza y lo que se necesita en cada uno, con alertas de información faltante;
  \item la verificación de reportes y de solicitudes;
  \item el seguimiento de cada solicitud hasta la entrega, con un código de un solo uso que se valida sin conexión;
  \item el registro de donaciones y de inventario por centro de acopio;
  \item un plan de asignación optimizado de los recursos a las solicitudes, dividido en paquetes pequeños que la entidad revisa y aprueba de uno en uno;
  \item un recorrido por etapas de cada reporte, solicitud y donación, y un panel que muestra dónde se acumula el trabajo;
  \item un mapa de la ruta de cada entrega, que ven la persona que pidió la ayuda, quien dona (de forma aproximada), el equipo de entrega y la entidad (control logístico). En esta versión la posición del equipo es \emph{simulada};
  \item la carga de datos oficiales abiertos sobre eventos y daños en catástrofes;
  \item un generador de escenarios simulados de varios tipos de desastre.
\end{itemize}

Queda fuera de esta versión: la conexión automática con sistemas oficiales, las fuentes satelitales o de drones, las redes sociales, los modelos de aprendizaje supervisado, la mensajería por SMS, la verificación de identidad contra bases oficiales y cualquier uso operativo real (ver sección~\ref{sec:fuera}).

\subsection{Contexto del reto}
El reto plantea que, tras un evento catastrófico, hay cortes de comunicación, rutas bloqueadas y una demanda de ayuda que desborda a los equipos de respuesta. Las decisiones se toman con rapidez y con datos incompletos. Los antecedentes revisados muestran que la priorización multicriterio está bien estudiada~\cite{Sahin2023314,Chakraborty2016308}, pero casi no se modela la confiabilidad del dato dentro de la decisión~\cite{Suarez2022443,Fan2020}. Existen también aplicaciones de evaluación de necesidades~\cite{Montefalcon202154} y sistemas nacionales como el brasileño S2iD~\cite{Eckhardt2018}, que sirven de referencia para comparar.

\subsection{Definiciones y siglas}
\begin{tabularx}{\textwidth}{@{}lX@{}}
\toprule
\textbf{Término} & \textbf{Significado} \\
\midrule
Municipio & Unidad administrativa del Atlántico (23 en total). \\
Barrio & Unidad espacial de análisis dentro de un municipio. En los municipios donde no hay barrios mapeados se usan los centros poblados. \\
Reporte & Observación enviada por una persona sobre el estado de un barrio. \\
Prioridad ($P$) & Puntaje de 0 a 1 que ordena los barrios según su necesidad de atención. \\
Confianza ($C$) & Valor de 0 a 1 que indica cuánto se puede confiar en la información de un barrio. \\
TOPSIS & Método multicriterio que ordena alternativas por su cercanía a la mejor solución. \\
AHP & Proceso analítico jerárquico, usado para obtener los pesos de los criterios. \\
MoSCoW & Priorización de requisitos: \Must{} (debe), \Should{} (debería), \Could{} (podría), \Wont{} (no en esta versión). \\
WCAG & Pautas de accesibilidad para contenido web. \\
RF / RNF & Requisito funcional / requisito no funcional. \\
Donación & Bien (o monto, si es monetaria) que una persona, organización o empresa ofrece para la respuesta. \\
Lote & Cantidad de un mismo bien, con el mismo estado y vencimiento, guardada en un centro de acopio. \\
Plan de asignación & Propuesta de qué lote se entrega a qué solicitud y en qué cantidad. \\
Paquete & Parte pequeña de un plan (hasta 6 entregas de un mismo municipio, con un solo equipo) que la entidad revisa y aprueba de una vez. \\
Recorrido & Las etapas por las que pasa un reporte, una solicitud o una donación, y en cuál va. \\
FEFO & ``Primero en vencer, primero en salir'': se entregan antes los lotes que vencen antes. \\
\bottomrule
\end{tabularx}
